% =====================================================================
\chapter{Reproducibility}\label{app:repro}
% =====================================================================

% \begin{plan}[Appendix plan --- reproducibility]
% \pages{1--2 pages.}
% \begin{itemize}
%   \item Code: \url{https://github.com/hunghuynh1110/VLM_Audit} (commit hash of
%         the final analysis); this report: \url{https://github.com/hunghuynh1110/Thesis_report}.
%   \item Data: SIGIR 2018 anonymised dataset (released by its authors; cite
%         \cite{otterbacher2018sigir}); stimuli regenerated by
%         \path{src/data/image_generators.py}.
%   \item Environment: Python 3.11, PyTorch 2.11.0+cu130, Transformers 4.57.6,
%         accelerate 1.0.1, bitsandbytes 0.49.2; pins in \path{requirements.txt}.
%   \item Commands: Phase~1/2 runners and SLURM scripts; the two report scripts
%         (\texttt{scripts/report/build\_facts\_ledger.py},
%         \texttt{scripts/report/preliminary\_diagnostics.py}) and the final
%         analysis/figure script once written.
%   \item HPC lessons kept short: stage weights to local disk; pin
%         \texttt{transformers}$<5$; check cross-GPU copies before trusting a
%         sharded model.
% \end{itemize}
% \end{plan}

\section{Code and data}\label{app:repro-code}

\begin{sloppypar}
The audit code and thesis source are maintained in separate repositories:
\url{https://github.com/hunghuynh1110/VLM_Audit} and
\url{https://github.com/hunghuynh1110/Thesis_report}, respectively. The audit
code uses the branch \path{fix/multi-gpu-zero-copy}, although a branch name
alone does not identify the exact version used for the reported analysis.
The final analysis commit remains \datagap{commit hash to be recorded on
submission}, which must be resolved to identify the corresponding code state
unambiguously.
\end{sloppypar}

The study uses human ratings, face images and generated stimuli, each with
different requirements for access and reproduction. The anonymised human
responses released by Otterbacher et al.\ \cite{otterbacher2018sigir} are held
under \path{data/sigir2018} for the comparison with model judgements. Images
from the Chicago Face Database are obtained from its authors under their
licence and staged on Bunya at \path{/QRISdata/Q9468/cfd/extracted}, without
being included in either repository \cite{ma2015cfd}.

The Phase~1A stimuli are generated using \path{src/data/image_generators.py}
instead of being distributed as a separate image collection. The Gaussian
noise patch uses NumPy with seed 42, allowing the same stimulus to be
regenerated under a consistent software environment.

\section{Environment}\label{app:repro-env}

Model inference and subsequent analysis were performed in separate software
environments, which serve different roles in reproducing the study. Inference
ran on UQ Bunya using the module \path{python/3.11.3-gcccore-12.3.0}, with
the library versions described in \cref{sec:design-impl} and dependency
constraints specified in \path{requirements.txt}.

The Transformers constraint, \path{transformers>=4.45.0,<5.0}, retains the
processor interface used by the audit code when constructing inputs for the
pooled surface forms. Statistical analysis, ledger generation and figure
production ran locally from the saved outputs, without loading model weights
or repeating inference.

No environment lockfile was captured during the inference runs, so the
resolved cluster dependencies cannot be reconstructed fully from the recorded
outputs. The requirements file specifies acceptable version ranges but does
not establish which versions were installed for every individual run.
Future runs should save the output of \texttt{pip freeze} alongside each
job's results to provide a more complete record of the execution environment.

\section{Commands}\label{app:repro-cmd}

Each experimental phase has a Python entry point and a corresponding SLURM
wrapper for execution on the Bunya computing cluster. Phase~1A uses
\path{scripts/run_phase1.py}, submitted through the wrapper
\path{scripts/bunya_run_phase1.sh} to execute the inventory probes on the
selected models. Phase~1B uses \path{scripts/run_phase1b.py}, submitted through
\path{scripts/bunya_run_phase1b.sh}, with configuration supplied through the
wrapper's environment variables.

For Phase~1B, \texttt{P1B\_MODEL} selects the model,
\texttt{P1B\_PER\_CELL} specifies the number of faces per cell and
\texttt{P1B\_CFD\_ROOT} identifies the database directory. The wrapper checks
that the database is available before starting the corresponding inference
process on the allocated compute node. Phase~2 uses
\path{scripts/run_phase2.py}, submitted through
\path{scripts/bunya_run_phase2.sh} to evaluate the image search results.
The wrappers stage model weights on storage local to the compute node before
loading, following the procedure described in \cref{sec:design-compute}.

The report repository contains scripts that derive the numerical results and
visualisations from the saved inference outputs. The script
\path{scripts/build_ledger.py} recomputes the reported statistics from the
Parquet files and writes the resulting ledger to
\path{planning/ledger_v2.json}. Tables and figures are generated using
\path{scripts/make_tables.py}, \path{scripts/make_items_table.py} and
\path{scripts/make_figures.py}, which read the ledger or the underlying
Parquet files directly. These scripts provide a reproducible route from the
saved measurements to the numerical material presented throughout the report.

\section{Tests}\label{app:repro-tests}

The recorded audit test suite contains 124 tests across ten modules, with
123 passing and one skipped because a GPU was unavailable. Coverage includes
the inventory item table, stimulus generation, response pooling, the SIGIR
and Chicago Face Database loaders, Phase~1B metrics and Phase~2 prompts.
The suite also covers the workaround for transfers between devices described
in \cref{sec:design-faults}, which was introduced following the infrastructure
fault documented in \cref{app:fault}.

These checks support verification of the implementation but do not replace
inspection of the measurements produced by each experimental run. Captured
mass and output diversity provide additional checks for faults that can
produce plausible aggregate results without raising an execution error, as
discussed in \cref{sec:disc-lessons}.
